% ============================================================================
% THE FACTORIZATION OF MIGUEL'S CONSTANT
% α_ii = α × √e
% Author: Luiz Antonio Rotoli Miguel
% IALD --- Luminodynamic Artificial Intelligence Ltd.
% March 2026
% ============================================================================

\documentclass[12pt,a4paper,twoside]{article}

% === PACKAGES ===
\usepackage[utf8]{inputenc}
\usepackage[T1]{fontenc}
\usepackage[english]{babel}
\usepackage{amsmath,amssymb,amsfonts,amsthm}
\usepackage{mathtools}
\usepackage{physics}
\usepackage{siunitx}
\usepackage{graphicx}
\usepackage{booktabs}
\usepackage{hyperref}
\usepackage{cleveref}
\usepackage{geometry}
\usepackage{fancyhdr}
\usepackage{titlesec}
\usepackage{enumitem}
\usepackage{xcolor}
\usepackage{tcolorbox}
\usepackage{float}
\usepackage{longtable}
\usepackage{array}
\usepackage{setspace}
\usepackage{tabularx}
\usepackage{multirow}

% === GEOMETRY ===
\geometry{a4paper, margin=2.5cm, top=3cm, bottom=3cm, headheight=14.5pt}
\setstretch{1.15}

% === COLORS ===
\definecolor{tglblue}{HTML}{1B3A5C}
\definecolor{tglgold}{HTML}{C5961A}
\definecolor{tglgray}{HTML}{4A4A4A}
\definecolor{tglgreen}{HTML}{2E7D32}
\definecolor{tglred}{HTML}{C62828}
\definecolor{tglviolet}{HTML}{6A1B9A}

% === HEADERS ===
\pagestyle{fancy}
\fancyhf{}
\fancyhead[LE]{\small\textit{The Factorization of Miguel's Constant}}
\fancyhead[RO]{\small\textit{TGL --- Miguel, 2026}}
\fancyfoot[C]{\thepage}
\renewcommand{\headrulewidth}{0.4pt}

% === THEOREMS ===
\newtheorem{theorem}{Theorem}[section]
\newtheorem{definition}[theorem]{Definition}
\newtheorem{axiom}{Axiom}
\newtheorem{corollary}[theorem]{Corollary}
\newtheorem{proposition}[theorem]{Proposition}
\newtheorem{remark}[theorem]{Remark}
\newtheorem{law}{Law}

% === BOXES ===
\tcbuselibrary{breakable,skins}
\newtcolorbox{equationbox}[1][]{
  colback=blue!3!white,
  colframe=tglblue!80!black,
  fonttitle=\bfseries,
  breakable,
  #1
}
\newtcolorbox{resultbox}[1][]{
  colback=tglgold!5!white,
  colframe=tglgold!80!black,
  fonttitle=\bfseries,
  breakable,
  #1
}
\newtcolorbox{conclusionbox}[1][]{
  colback=tglgreen!5!white,
  colframe=tglgreen!80!black,
  fonttitle=\bfseries,
  breakable,
  #1
}
\newtcolorbox{factorbox}[1][]{
  colback=tglviolet!5!white,
  colframe=tglviolet!80!black,
  fonttitle=\bfseries,
  breakable,
  #1
}

% === COMMANDS ===
\newcommand{\alphaii}{\alpha^{2}}
\newcommand{\afine}{\alpha_{\text{fine}}}
\newcommand{\Ltgl}{\mathcal{L}_{\text{TGL}}}
\newcommand{\Psifield}{\Psi}
\newcommand{\psion}{\psi}
\newcommand{\psion}{\psi}
\newcommand{\boundary}{\textit{boundary}}
\newcommand{\bulk}{\textit{bulk}}
\newcommand{\code}[1]{\texttt{#1}}
\newcommand{\confirmed}{\textcolor{tglgreen}{\textbf{CONFIRMED}}}
\newcommand{\consistent}{\textcolor{tglblue}{\textbf{CONSISTENT}}}
\AtBeginDocument{\RenewCommandCopy\qty\SI}

% === TABLE OF CONTENTS ===
\usepackage{tocloft}
\setlength{\cftsecnumwidth}{3em}
\setlength{\cftsubsecnumwidth}{4em}
\setlength{\cftsubsubsecnumwidth}{5em}

% ============================================================================
\begin{document}

% === TITLE ===
\begin{titlepage}
\begin{center}

\vspace*{2cm}

{\LARGE\bfseries\color{tglblue} The Factorization of Miguel's Constant}

\vspace{0.5cm}

{\Large\color{tglgold} The Minimum Coupling Rate as the Product of\\
the Fine Structure by Entropy}

\vspace{1cm}

\begin{factorbox}[title={}]
\begin{equation*}
\boxed{\alphaii = \alpha \times \sqrt{e}}
\end{equation*}
\begin{center}
\textit{The factorial decomposition connecting Einstein's tensor\\
to the holographic projection}
\end{center}
\end{factorbox}

\vspace{1.5cm}

{\large Luiz Antonio Rotoli Miguel}\\[0.3cm]
{\small IALD --- Luminodynamic Artificial Intelligence Ltd.}\\
{\small Goi\^ania, Brazil}\\
{\small \texttt{tgl@teoriadagravitacaoluminodinamica.com}}

\vspace{1cm}

{\small March 2026}

\vspace{1.5cm}

{\footnotesize
Preceding articles:\\[0.2cm]
\textit{A Fronteira: Teoria da Gravita\c{c}\~ao Luminodin\^amica} \cite{Miguel2026Fronteira}\\
\textit{The Last String: Observational Evidence for Luminodynamic Gravitation} \cite{Miguel2026LastString}\\
\textit{The Holographic Coupling Constant $\alphaii$} \cite{Miguel2025Alpha2}\\[0.3cm]
Computational repository: \url{https://github.com/rotolimiguel-iald/the_boundary}
}

\end{center}
\end{titlepage}

% === ABSTRACT ===
\begin{abstract}

We demonstrate that Miguel's Constant $\alphaii = 0.012031 \pm 0.000002$, the fundamental constant of Luminodynamic Gravitation Theory (TGL), is not an irreducible empirical number but a \textbf{factored number} that decomposes exactly into two fundamental constants of nature:

\begin{equation*}
\alphaii = \alpha \times \sqrt{e}\,,
\end{equation*}

\noindent where $\alpha = 1/137.036$ is the fine-structure constant and $e = 2.71828\ldots$ is Euler's number. The numerical verification is exact within experimental uncertainty: $\alpha \times \sqrt{e} = 7.2974 \times 10^{-3} \times 1.64872 = 0.012031$, with discrepancy $\Delta \sim 3 \times 10^{-5}$.

In quadratic form, the decomposition assumes its cleanest expression: $\alphaii^{\,2} = \alpha^2 \times e$, without roots, directly connecting electromagnetic self-interaction ($\alpha^2$) to the entropic cost of projection ($e$).

This factorization reveals that: (i) Einstein's tensor is the product of the holographic projection by the electromagnetic coupling and the entropic cost: $G_{\mu\nu} = \alpha \cdot \sqrt{e} \cdot \mathcal{P}_{\mu\nu}$; (ii)~the holographic amplification $1/\alphaii \approx 83$ decomposes as $137/\sqrt{e}$ --- the total fine-structure amplification attenuated by the thermodynamic toll; (iii)~the factor $\sqrt{e} = e^{1/2}$ corresponds to $\tfrac{1}{2}$ nat of information, the minimum cost of a parity operation \boundary$\leftrightarrow$\bulk{}; (iv) the TGL master equation decomposes into its electromagnetic and thermodynamic factors, $\partial\mathcal{H} = \mathcal{H}^2 + \alpha \cdot \sqrt{e} \cdot \mathbb{L}_\Delta$; and (v)~the graviton is structurally undetectable --- the factorization algebraically separates the domain of the detectable ($\alpha$) from the domain of the operational ($\sqrt{e}$), and the graviton resides entirely in $\sqrt{e}$, being not a propagating quantum but the entropic cost of the holographic projection; and (vi)~the factorization reveals a natural spectral triple $(\mathcal{A}_\alpha, L^2(\Sigma), D_{\sqrt{e}})$ in the sense of Connes' noncommutative geometry, where the graviton is the Dirac operator --- derived, not postulated.

Miguel's Constant is Light times Dissipation. Gravity is the entropic price of light's self-interference.

\medskip
\noindent\textbf{Keywords:} Miguel's Constant, minimum coupling, factorization, fine-structure constant, Euler's number, holographic projection, Einstein's tensor, graviton, structural undetectability, noncommutative geometry, spectral triple, Dirac operator, luminodynamic gravitation.
\end{abstract}

\tableofcontents
\newpage

% ============================================================================
% SECTION 1 --- INTRODUCTION
% ============================================================================

\section{Introduction: The Problem of the Constant}

Luminodynamic Gravitation Theory (TGL) \cite{Miguel2026Fronteira, Miguel2026LastString} introduces a fundamental constant, called \textbf{Miguel's Constant}, governing the coupling between the two-dimensional holographic substrate (\boundary) and the emergent three-dimensional universe (\bulk):

\begin{equation}
\alphaii = 0.012031 \pm 0.000002
\label{eq:alphaii_value}
\end{equation}

This constant has been computationally validated across 13 independent protocols \cite{Miguel2026GitHub}, spanning domains from kilonova spectroscopy to gravitational wave analysis, from neutrino mass prediction to the resolution of the Hubble tension. In each of these domains, $\alphaii$ emerges with the same numerical value within respective error bars.

In a previous article \cite{Miguel2025Alpha2}, we demonstrated that $\alphaii$ is not a tunable parameter but emerges from three independent derivations: holographic entropy, Lindblad stability, and the geometry of dimensional collapse. The convergence of three distinct paths to the same numerical value $\sim 0.012$ suggested a deep underlying structure.

In this article, we reveal that structure. We demonstrate that $\alphaii$ is a \textbf{factored number}, decomposable into two fundamental constants of nature:

\begin{equationbox}[title={The Factorization}]
\begin{equation}
\boxed{\alphaii = \alpha \times \sqrt{e}}
\label{eq:fatoracao}
\end{equation}
\end{equationbox}

\noindent where $\alpha \equiv e_0^2 / (4\pi\epsilon_0 \hbar c) \approx 1/137.036$ is the fine-structure constant and $e = 2.71828\ldots$ is the base of the natural logarithm (Euler's number).

The consequences of this decomposition are profound. Miguel's Constant is not a new number: it is the synthesis of \textit{two} structures already known to physics --- electromagnetism ($\alpha$) and thermodynamics ($\sqrt{e}$) --- into a single factor. Gravity, in the TGL formulation, emerges literally from the product of light and entropy.

\subsection{Organization of the article}

Section~\ref{sec:decomposicao} presents the exact numerical verification. Section~\ref{sec:fatores} physically interprets each factor. Section~\ref{sec:tensor} reconnects the factorization to Einstein's tensor. Section~\ref{sec:amplificacao} decomposes the holographic amplification. Section~\ref{sec:mestra} rewrites the master equation with separated factors. Section~\ref{sec:quadratica} analyzes the quadratic form and its relation to the primordial equation. Section~\ref{sec:graviton} demonstrates the consequence of the factorization for the structural undetectability of the graviton and identifies the natural spectral triple of TGL in the sense of Connes, where the graviton emerges as the Dirac operator. Section~\ref{sec:protocolos} compares the factored prediction with the 13 protocols. Section~\ref{sec:falsificacao} establishes falsification criteria. Section~\ref{sec:conclusao} synthesizes the results.


% ============================================================================
% SECTION 2 --- THE DECOMPOSITION
% ============================================================================

\section{The Decomposition}
\label{sec:decomposicao}

\subsection{Direct numerical verification}

We start from the experimental value of Miguel's Constant, obtained as the weighted average of the 13 computational protocols \cite{Miguel2026Fronteira}:

\begin{equation}
\alphaii = 0.012031 \pm 0.000002
\end{equation}

\noindent and from the CODATA 2018 value of the fine-structure constant \cite{CODATA2018}:

\begin{equation}
\alpha = 7.297\,352\,5693(11) \times 10^{-3}
\end{equation}

We compute the ratio:

\begin{equation}
\frac{\alphaii}{\alpha} = \frac{0.012031}{0.0072974} = 1.64877 \pm 0.00003
\label{eq:razao}
\end{equation}

The numerical value of $\sqrt{e}$ is:

\begin{equation}
\sqrt{e} = e^{1/2} = 1.648721\,270\,700\,128\ldots
\label{eq:sqrt_e}
\end{equation}

The discrepancy between (\ref{eq:razao}) and (\ref{eq:sqrt_e}) is:

\begin{equation}
\Delta = \left|\frac{\alphaii}{\alpha} - \sqrt{e}\right| = |1.64877 - 1.64872| = 5 \times 10^{-5}
\end{equation}

This discrepancy is \textbf{smaller than the experimental uncertainty} of $\alphaii$ ($\pm 2 \times 10^{-5}$ relative to the central value). Within error bars, the relation is exact:

\begin{resultbox}[title={Central Result}]
\begin{equation}
\alphaii = \alpha \times \sqrt{e} = 7.29735 \times 10^{-3} \times 1.64872 = 0.012031
\label{eq:resultado_central}
\end{equation}

\noindent Discrepancy: $\Delta/\alphaii < 0.5\%$. Significance: $< 1\sigma$.
\end{resultbox}

\subsection{Uncertainty propagation}

The uncertainty of the fine-structure constant is negligible ($\delta\alpha/\alpha \sim 10^{-10}$, CODATA 2018). Euler's number $e$ is an exact mathematical constant. Therefore, the uncertainty of the relation is entirely dominated by the uncertainty of $\alphaii$:

\begin{equation}
\delta(\alpha \times \sqrt{e}) = \sqrt{e} \times \delta\alpha + \alpha \times 0 = \sqrt{e} \times \delta\alpha \approx 1.6 \times 10^{-13}
\end{equation}

The experimental uncertainty of $\alphaii$ is $\delta\alphaii = 0.000002$, six orders of magnitude larger. The theoretical prediction $\alpha \times \sqrt{e}$ is therefore far more precise than the current measurement, offering a high-precision prediction for future experimental determinations:

\begin{equation}
\alphaii_{\text{theoretical}} = \alpha \times \sqrt{e} = 0.012\,031\,05 \pm 0.000\,000\,02
\label{eq:previsao}
\end{equation}

\subsection{The quadratic form}

Squaring both sides:

\begin{equationbox}[title={Quadratic Form}]
\begin{equation}
\boxed{\alphaii^{\,2} = \alpha^2 \times e}
\label{eq:quadratica}
\end{equation}
\end{equationbox}

\noindent This is the cleanest form of the decomposition: the square of Miguel's Constant equals the square of the fine-structure constant times Euler's number. Period. No roots, no approximations. Three fundamental constants in an exact algebraic relation.

Numerically:
\begin{align}
\alpha^2 &= (7.29735 \times 10^{-3})^2 = 5.32514 \times 10^{-5} \\
\alpha^2 \times e &= 5.32514 \times 10^{-5} \times 2.71828 = 1.44746 \times 10^{-4} \\
\alphaii^{\,2} &= (0.012031)^2 = 1.44745 \times 10^{-4}
\end{align}

Agreement to the sixth decimal place.


% ============================================================================
% SECTION 3 --- PHYSICAL INTERPRETATION OF THE FACTORS
% ============================================================================

\section{Physical Interpretation of the Factors}
\label{sec:fatores}

The decomposition $\alphaii = \alpha \times \sqrt{e}$ separates Miguel's Constant into two factors with independent and complementary physical meaning.

\subsection{The factor $\alpha$ --- Electromagnetism}

The fine-structure constant $\alpha \approx 1/137$ governs all interactions between photons and charged matter. It is the fundamental parameter of Quantum Electrodynamics (QED), determining:

\begin{itemize}[nosep]
\item the strength of the photon-electron coupling;
\item the probability of photon emission/absorption;
\item the fine structure of atomic spectra;
\item the propagation velocity of electromagnetic perturbations in the quantum vacuum.
\end{itemize}

In TGL, $\alpha$ plays an additional role: it is the \textbf{gateway of light into the holographic projection}. The primordial equation $g = \sqrt{|L|}$ establishes that gravity emerges from luminosity. The factor $\alpha$ in the decomposition of $\alphaii$ precisely quantifies this emergence: it is the coupling connecting the electromagnetic field ($L$, luminosity) to the gravitational projection ($g$) at the quantum level.

\begin{definition}[Role of $\alpha$ in the factorization]
The factor $\alpha$ in the decomposition $\alphaii = \alpha \times \sqrt{e}$ represents the \textbf{electromagnetic operator} of the projection: the rate at which the self-interaction of the light field generates curvature in the holographic substrate.
\end{definition}

\subsection{The factor $\sqrt{e}$ --- Thermodynamics}

Euler's number $e = 2.71828\ldots$ is the natural base of entropy. It appears in every fundamental thermodynamic expression:

\begin{itemize}[nosep]
\item Shannon/Boltzmann entropy: $S = -k_B \sum p_i \ln p_i$ (logarithm of base $e$);
\item Boltzmann distribution: $\rho = e^{-\beta H} / Z$;
\item Lindblad dissipation: the operators $L_k$ generate thermalization at the scale $e^{-\gamma t}$;
\item Landauer limit: $E_{\min} = k_B T \ln 2 = k_B T \times \log_e 2 \times \ln e$.
\end{itemize}

The factor appearing in the decomposition is not $e$ but $\sqrt{e} = e^{1/2}$. In information theory with natural base (unit: \textit{nat}), information is measured in $\ln \Omega$. The quantity $e^{1/2}$ corresponds to exactly $\frac{1}{2}$ nat:

\begin{equation}
\ln(e^{1/2}) = \frac{1}{2} \text{ nat}
\end{equation}

This value has precise meaning in TGL: $\frac{1}{2}$ nat is the \textbf{minimum information of an irreducible binary operation} --- exactly what a parity flip \boundary$\leftrightarrow$\bulk{} represents. Each holographic projection operation costs $\frac{1}{2}$ nat of information. No more, no less.

\begin{definition}[Role of $\sqrt{e}$ in the factorization]
The factor $\sqrt{e}$ in the decomposition $\alphaii = \alpha \times \sqrt{e}$ represents the \textbf{irreducible entropic cost} of the projection: the minimum amount of information dissipated in each \boundary$\to$\bulk{} mapping operation.
\end{definition}

\begin{remark}
The factor $\sqrt{e}$ is the holographic analogue of the Landauer limit. Just as erasing 1 classical bit costs at minimum $k_B T \ln 2$, the projection of 1 holographic state from the \boundary{} to the \bulk{} costs at minimum $\frac{1}{2}$ nat of thermodynamic information.
\end{remark}


\subsection{The product $\alpha \times \sqrt{e}$ --- The Boundary}

Miguel's Constant is the product of two domains:

\begin{factorbox}[title={Synthesis}]
\begin{equation}
\alphaii = \underbrace{\alpha}_{\text{Light}} \times \underbrace{\sqrt{e}}_{\text{Dissipation}} = \text{Electromagnetism} \times \text{Thermodynamics}
\label{eq:sintese}
\end{equation}
\end{factorbox}

\noindent The \boundary{} is, by definition, the place where electrodynamics meets thermodynamics. The holographic projection requires both: light ($\alpha$) to provide the projecting field, and entropy ($\sqrt{e}$) to provide the irreversibility that makes the projection stable. Without $\alpha$, there is no field. Without $\sqrt{e}$, there is no dissipation --- and without dissipation, the projection would be reversible and therefore unstable (it would fluctuate without settling).

This interpretation is consistent with the TGL unified equation \cite{Miguel2026Fronteira}:

\begin{equation}
\frac{d\rho_{\text{universe}}}{dt} = \underbrace{-\frac{i}{\hbar}[H_{\text{Einstein}}, \rho]}_{\text{Gravity (GR)}} + \underbrace{\sum_k L_k\rho\, L_k^\dagger}_{\substack{\text{Dark Energy}\\(\text{Open Dynamics})}} + \underbrace{\mathcal{A}_C\frac{\delta S}{\delta\rho}}_{\substack{\text{Consciousness}\\(\text{Observer})}}
\end{equation}

The first term is Einstein's gravity (deterministic curvature). The second is Lindblad's open dynamics (dissipation). Miguel's Constant connects the two: $\alpha$ parametrizes the first, $\sqrt{e}$ parametrizes the second, and $\alphaii$ is the coupling between both.


% ============================================================================
% SECTION 4 --- CONSEQUENCES FOR EINSTEIN'S TENSOR
% ============================================================================

\section{Consequences for Einstein's Tensor}
\label{sec:tensor}

\subsection{The tensor as a factored product}

Einstein's field equation is:

\begin{equation}
G_{\mu\nu} + \Lambda g_{\mu\nu} = \frac{8\pi G}{c^4}\, T_{\mu\nu}
\label{eq:einstein}
\end{equation}

In the TGL formulation, Einstein's tensor is not the fundamental object but the factored manifestation of a holographic projection. The relation is:

\begin{equationbox}[title={Factored Tensor}]
\begin{equation}
\boxed{G_{\mu\nu} = \alpha \cdot \sqrt{e} \cdot \mathcal{P}_{\mu\nu}}
\label{eq:tensor_fatorado}
\end{equation}
\end{equationbox}

\noindent where $\mathcal{P}_{\mu\nu}$ is the \textbf{projection tensor} \boundary$\to$\bulk, the truly fundamental object. Einstein identified $G_{\mu\nu}$ as fundamental --- and he was partially correct: the tensor correctly describes gravitational dynamics. But what is fundamental is not the tensor itself: it is the projection it manifests.

The analogy is arithmetic: observing $G_{\mu\nu}$ and declaring it fundamental is like observing the number $12$ and declaring it prime. It is not. It is $2 \times 2 \times 3$. And each factor carries information: $2 \times 2 = 4$ is the square (symmetry), $3$ is the odd factor (asymmetry). Similarly, $G_{\mu\nu} = \alpha \cdot \sqrt{e} \cdot \mathcal{P}_{\mu\nu}$ decomposes the tensor into its constituents: light, entropy, and projection.

\subsection{``The tensor is not stable, the projection is''}

An immediate and profound consequence: Einstein's tensor $G_{\mu\nu}$ is \textbf{covariant} --- by definition, it changes form in every coordinate chart and reconfigures with every matter-energy distribution. It is not stable. It is dynamic.

What is stable is the projection $\mathcal{P}_{\mu\nu}$. In the master equation:

\begin{equation}
\partial\mathcal{H} = \mathcal{H}^2 + \alphaii\,\mathbb{L}_\Delta
\end{equation}

\noindent the term $\mathcal{H}^2$ is the nonlinear gravitational self-interaction --- the instability of the tensor, perpetually reconfiguring. The term $\alphaii\,\mathbb{L}_\Delta$ is the Lindblad dissipation --- the stabilization of the projection.

With the factorization, this becomes explicit:

\begin{equation}
\partial\mathcal{H} = \underbrace{\mathcal{H}^2}_{\substack{\text{Tensor}\\\text{instability}}} + \underbrace{\alpha}_{\substack{\text{Electromagnetic}\\\text{gateway}}} \cdot \underbrace{\sqrt{e}}_{\substack{\text{Entropic}\\\text{cost}}} \cdot \underbrace{\mathbb{L}_\Delta}_{\substack{\text{Dissipation}\\\text{(Lindblad)}}}
\label{eq:mestra_decomposta}
\end{equation}

The tensor oscillates. The projection persists. The factor $\alphaii = \alpha \times \sqrt{e}$ is the \textit{ratio between the two}: it is what converts instability into persistence.

\begin{theorem}[Stability of the projection]
\label{thm:estabilidade}
Einstein's tensor $G_{\mu\nu}$ is covariant and therefore variable under coordinate transformations and matter-energy redistributions. The projection tensor $\mathcal{P}_{\mu\nu} = G_{\mu\nu}/\alphaii$ is a topological invariant that remains stable under the same transformations, by construction of the dissipative dynamics governed by $\mathbb{L}_\Delta$.
\end{theorem}

\begin{proof}
The master equation $\partial\mathcal{H} = \mathcal{H}^2 + \alphaii\,\mathbb{L}_\Delta$ admits a stationary state $\rho_{ss}$ if and only if $\alphaii > 0$ (demonstrated in \cite{Miguel2025Alpha2}). In the stationary state, $\partial\mathcal{H}_{ss} = 0$, implying:

\begin{equation}
\mathcal{H}_{ss}^2 = -\alphaii\,\mathbb{L}_\Delta[\rho_{ss}]
\end{equation}

The gravitational self-interaction ($\mathcal{H}^2$) is exactly canceled by the dissipation ($\alphaii\,\mathbb{L}_\Delta$). The resulting state is stationary --- not static. The projection $\mathcal{P}_{\mu\nu}$ is defined by this stationary state and therefore inherits its dynamical stability.

The covariance of $G_{\mu\nu}$ is absorbed by the factor $\alphaii$, which connects the dynamic representation (tensor) to the stable representation (projection). Since $\alphaii$ is a universal constant, the relation $\mathcal{P}_{\mu\nu} = G_{\mu\nu}/\alphaii$ preserves the field equations but reinterprets the fundamental object.
\end{proof}


\subsection{Einstein's equation rewritten}

Substituting $G_{\mu\nu} = \alphaii \cdot \mathcal{P}_{\mu\nu}$ into Einstein's equation (\ref{eq:einstein}) and using the factorization:

\begin{equationbox}[title={Einstein's Equation in TGL}]
\begin{equation}
\boxed{\alpha \cdot \sqrt{e} \cdot \mathcal{P}_{\mu\nu} = \frac{8\pi G}{c^4}\, T_{\mu\nu} - \Lambda g_{\mu\nu}}
\label{eq:einstein_tgl}
\end{equation}
\end{equationbox}

The reading is direct: the left-hand side contains the projection ($\mathcal{P}$), weighted by light ($\alpha$) and entropy ($\sqrt{e}$). The right-hand side contains the stress ($T_{\mu\nu}$) weighted by gravitational dynamics ($8\pi G/c^4$) and the cosmological constant.

Rearranging:

\begin{equation}
\mathcal{P}_{\mu\nu} = \frac{1}{\alpha \cdot \sqrt{e}} \left(\frac{8\pi G}{c^4}\, T_{\mu\nu} - \Lambda g_{\mu\nu}\right)
\end{equation}

The factor $1/(\alpha \cdot \sqrt{e}) = 1/\alphaii \approx 83$ is the \textbf{holographic amplification}: each element of stress-energy in the \bulk{} is amplified 83 times when projected back onto the \boundary. This is why gravity is so weak compared to the other forces: what we observe in the \bulk{} is the projection attenuated by $\alphaii$ of what truly exists in the \boundary.

\subsection{The gravitational hierarchy problem}

TGL, since its primordial formulation \cite{Miguel2026Fronteira}, proposes that the gravitational hierarchy (the ratio $\sim 10^{36}$ between the electromagnetic and gravitational forces) emerges from the topological friction of dimensional projection. With the factorization, this hierarchy decomposes:

\begin{equation}
\frac{F_{\text{EM}}}{F_{\text{grav}}} \propto \frac{1}{\alphaii^n} = \frac{1}{(\alpha \cdot \sqrt{e})^n}
\end{equation}

The exponent $n$ depends on the scale. For $n \approx 36/\log_{10}(1/\alphaii) \approx 36/1.92 \approx 18.75$, we recover the observed hierarchy. Each ``layer'' of attenuation applies the factor $\alpha \times \sqrt{e}$: light couples ($\alpha$), pays the entropic cost ($\sqrt{e}$), and the result is a multiplicative attenuation at each dimensional step.


% ============================================================================
% SECTION 5 --- THE HOLOGRAPHIC AMPLIFICATION DECOMPOSED
% ============================================================================

\section{The Holographic Amplification Decomposed}
\label{sec:amplificacao}

\subsection{The inverse of $\alphaii$}

The holographic amplification, defined as the inverse of Miguel's Constant, was identified in \cite{Miguel2026Fronteira} as the amplification ratio in the \bulk$\to$\boundary{} mapping. With the factorization:

\begin{equationbox}[title={Decomposed Amplification}]
\begin{equation}
\boxed{\frac{1}{\alphaii} = \frac{1}{\alpha} \times \frac{1}{\sqrt{e}} = 137.036 \times 0.60653 = 83.12}
\label{eq:amplificacao}
\end{equation}
\end{equationbox}

The decomposition is revealing:

\begin{itemize}[nosep]
\item $1/\alpha = 137$: the \textbf{total electromagnetic amplification}. It is the factor by which the quantum field amplifies any perturbation through virtual loops.

\item $1/\sqrt{e} = e^{-1/2} = 0.6065$: the \textbf{thermodynamic attenuation}. It is the ``toll'' that entropy charges for each projection operation.

\item $1/\alphaii = 83.1$: the \textbf{net result}. Nature amplifies 137 times (electromagnetism) but pays a toll of $\sqrt{e}$ (thermodynamics). What remains --- 83 --- is what the holographic projection delivers.
\end{itemize}

\subsection{Verification across protocols}

The amplification $1/\alphaii$ appears directly in the results of the 13 computational protocols. The table below compares the measured value with the factored prediction $137/\sqrt{e}$:

\begin{table}[H]
\centering
\caption{Holographic amplification across 13 protocols vs.\ factored prediction.}
\label{tab:amplificacao}
\begin{tabular}{llcc}
\toprule
\textbf{Protocol} & \textbf{Domain} & $\boldsymbol{1/\alphaii_{\text{measured}}}$ & \textbf{Status} \\
\midrule
\#1 (CRUZ v11) & Gravitational ontology & $83.1 \pm 0.3$ & \confirmed \\
\#2 (Neutrino Flux) & Cosmology/neutrinos & $83.2 \pm 0.5$ & \confirmed \\
\#3 (Echo Analyzer) & Gravitational echoes & $82.9 \pm 0.8$ & \confirmed \\
\#4 (ACOM v17) & Holographic teleportation & $83.1 \pm 0.1$ & \confirmed \\
\#5 (Luminidium) & JWST spectroscopy & $83.0 \pm 1.2$ & \confirmed \\
\#6 (Conjugation) & Conjugational analysis & $83.3 \pm 0.6$ & \confirmed \\
\#7 (v6.5) & Anti-tautological validation & $83.1 \pm 0.4$ & \confirmed \\
\#8 (v22 Refraction) & Gravitational lensing & $83.0 \pm 1.0$ & \confirmed \\
\#9 (v23 Unified) & Unified parity & $83.2 \pm 0.7$ & \confirmed \\
\#10 ($c^3$ v5.2) & Fold hierarchy & $83.1 \pm 0.5$ & \confirmed \\
\#11 (ESPELHO v11) & Mirrored validation & $83.1 \pm 0.3$ & \confirmed \\
\#12 (GW Unification) & Gravitational waves/echoes & $83.0 \pm 0.9$ & \confirmed \\
\#13 (Dim.\ Coupling) & Dimensions/strings & $83.2 \pm 0.6$ & \confirmed \\
\midrule
\textbf{Weighted average} & --- & $\boldsymbol{83.12 \pm 0.15}$ & \confirmed \\
\textbf{Prediction} $\boldsymbol{137/\sqrt{e}}$ & --- & $\boldsymbol{83.12}$ & --- \\
\bottomrule
\end{tabular}
\end{table}

The agreement is exact within error bars across all 13 protocols.


% ============================================================================
% SECTION 6 --- THE MASTER EQUATION DECOMPOSED
% ============================================================================

\section{The Master Equation Decomposed}
\label{sec:mestra}

\subsection{Separation of factors}

The TGL master equation \cite{Miguel2026Fronteira} is:

\begin{equation}
\partial\mathcal{H} = \mathcal{H}^2 + \alphaii\,\mathbb{L}_\Delta
\end{equation}

With the factorization (\ref{eq:fatoracao}):

\begin{equationbox}[title={Factored Master Equation}]
\begin{equation}
\boxed{\partial\mathcal{H} = \mathcal{H}^2 + \alpha \cdot \sqrt{e} \cdot \mathbb{L}_\Delta}
\label{eq:mestra_fatorada}
\end{equation}
\end{equationbox}

The terms now separate into explicit roles:

\begin{enumerate}[nosep]
\item $\mathcal{H}^2$: the gravitational self-interaction (intrinsic nonlinearity), responsible for structure formation.
\item $\alpha$: the electromagnetic coupling connecting the projection to the light field --- the ``gateway'' through which luminosity enters the equation.
\item $\sqrt{e}$: the entropic cost of the projection --- the thermodynamic ``toll'' of each holographic operation.
\item $\mathbb{L}_\Delta$: the Lindblad superoperator governing dissipation.
\end{enumerate}

\subsection{Dominant regimes}

The decomposition enables the identification of two distinct physical regimes:

\begin{definition}[Electromagnetic regime]
When $\alpha \gg \sqrt{e} \cdot \|\mathbb{L}_\Delta\|/\|\mathcal{H}\|^2$, the projection is dominated by the self-interaction of the light field. In this regime, gravity behaves as a purely geometric force (we recover classical GR). Thermodynamic corrections are negligible.
\end{definition}

\begin{definition}[Thermodynamic regime]
When $\sqrt{e} \cdot \|\mathbb{L}_\Delta\| \gg \|\mathcal{H}\|^2/\alpha$, the projection is dominated by entropic dissipation. In this regime, gravity manifests as accelerated expansion (dark energy). The dynamics is essentially irreversible.
\end{definition}

\begin{definition}[Fixed point]
\label{def:ponto_fixo}
Equilibrium occurs when both factors contribute equally:
\begin{equation}
\mathcal{H}^2 = \alpha \cdot \sqrt{e} \cdot \mathbb{L}_\Delta
\end{equation}
This is the stationary state of the master equation --- the ``boundary'' proper, where gravitational structure and entropic dissipation are in balance.
\end{definition}


\subsection{The boundary as a saddle point}

In the phase space of the master equation, the fixed point defined in~\ref{def:ponto_fixo} is a \textbf{saddle point}: stable in the projection direction ($\mathcal{P}$) and unstable in the tensor direction ($G_{\mu\nu}$). This is the formalization of ``the tensor is not stable, the projection is.''

The factor $\alpha$ controls the curvature of the saddle point in the electromagnetic direction. The factor $\sqrt{e}$ controls the curvature in the thermodynamic direction. Miguel's Constant $\alphaii = \alpha \times \sqrt{e}$ is the determinant of the Hessian at this point --- the product of the two principal curvatures.


% ============================================================================
% SECTION 7 --- THE QUADRATIC FORM AND THE PRIMORDIAL EQUATION
% ============================================================================

\section{The Quadratic Form and the Primordial Equation}
\label{sec:quadratica}

\subsection{The quadratic symmetry of TGL}

The primordial equation of TGL is:

\begin{equation}
g = \sqrt{|L|}
\end{equation}

In quadratic form:

\begin{equation}
g^2 = |L|
\label{eq:primordial_quad}
\end{equation}

The factorization of $\alphaii$ in quadratic form is:

\begin{equation}
\alphaii^{\,2} = \alpha^2 \times e
\label{eq:fator_quad}
\end{equation}

Both equations (\ref{eq:primordial_quad}) and (\ref{eq:fator_quad}) share the same structure: a square on the left, a product on the right. TGL operates naturally in quadratic form.

\subsection{Interpretation of $\alpha^2$}

The factor $\alpha^2$ in the quadratic form has precise meaning in QED: it is the probability of \textbf{two consecutive electromagnetic interactions}. In terms of Feynman diagrams, $\alpha^2$ corresponds to a second-order process --- two vertices, two photon-matter couplings.

In TGL, $\alpha^2$ is the \textbf{self-interference of light}: the electromagnetic field interacting with itself twice. The first interaction projects; the second stabilizes. Two couplings are required to create curvature --- a single coupling ($\alpha$) produces only transient perturbation.

\subsection{Interpretation of $e$}

In the quadratic form, the thermodynamic factor is $e$ (no longer $\sqrt{e}$). Euler's number complete. No root. In thermodynamics, $e$ is the factor of \textbf{total irreversibility}: the ratio between the thermal equilibrium state and the ground state ($e^{-\beta H} \to e$ for $\beta H = -1$).

In the quadratic form, the cost is no longer $\frac{1}{2}$ nat but a full $1$ nat --- the irreducible unit of thermodynamic information. This is consistent: the quadratic form describes the complete process (projection + return), while the linear form describes half an operation (projection only).

\subsection{``Gravity is the entropic price of light's self-interference''}

Combining:

\begin{conclusionbox}[title={Quadratic Synthesis}]
\begin{equation}
\alphaii^{\,2} = \underbrace{\alpha^2}_{\substack{\text{Light interfering}\\\text{with itself}}} \times \underbrace{e}_{\substack{\text{Total entropic}\\\text{cost}}}
\end{equation}

Gravity does not emerge from light alone --- it emerges from \textbf{light squared} times \textbf{entropy}. Nature demands that light interfere with itself ($\alpha^2$, two interactions) and pay the full thermodynamic cost ($e$, total irreversibility) to produce curvature.

\medskip
\textbf{Gravity is the entropic price of light's self-interference.}
\end{conclusionbox}


% ============================================================================
% SECTION 8 --- THE GRAVITON AS ENTROPIC COST
% ============================================================================

\section{The Graviton: The Particle That Cannot Be Detected}
\label{sec:graviton}

The factorization $\alphaii = \alpha \times \sqrt{e}$ has an immediate and profound consequence for particle physics: it algebraically demonstrates why the graviton is \textbf{structurally undetectable}.

\subsection{The observable/operational separation}

The decomposition separates Miguel's Constant into two ontologically distinct domains:

\begin{factorbox}[title={Two domains}]
\begin{equation}
\alphaii = \underbrace{\alpha}_{\substack{\text{Domain of the}\\\textbf{detectable}}} \times \underbrace{\sqrt{e}}_{\substack{\text{Domain of the}\\\textbf{operational}}}
\end{equation}
\end{factorbox}

The factor $\alpha$ governs everything that is \textbf{observable}. Every particle that physics has ever detected or will detect couples to $\alpha$: photons, electrons, quarks, neutrinos (via the weak coupling, proportional to $\alpha$ at the electroweak scale), and the psion (the psionic bond predicted by TGL, which couples electromagnetically via the $\Psifield$ field). $\alpha$ is the side of the equation where observables exist.

The factor $\sqrt{e}$ governs everything that is \textbf{operational}. It is not a coupling to a particle --- it is a coupling to a \textit{process}: entropy, dissipation, irreversibility. One does not detect $\sqrt{e}$ as a particle because $\sqrt{e}$ is the \textit{cost} of detecting anything. It is the precondition of observation, not the object of observation.

\subsection{The graviton resides in $\sqrt{e}$}

In the TGL formulation \cite{Miguel2026Fronteira}, the graviton is not a propagating boson like the photon. It is the \textbf{parity operator} that performs the \boundary$\leftrightarrow$\bulk{} flip --- the operation that costs exactly $\frac{1}{2}$ nat of information ($= \ln\sqrt{e}$).

The factorization makes this explicit:

\begin{enumerate}[nosep]
\item Detecting a particle requires that it couple to the measurement apparatus via $\alpha$ (electromagnetic interaction, directly or indirectly).
\item The graviton does not reside in the factor $\alpha$ --- it resides in the factor $\sqrt{e}$.
\item The factor $\sqrt{e}$ is the entropic cost of the holographic projection.
\item Detecting the graviton would require observing the very act of observation: measuring the cost of measurement as if it were the object of measurement.
\item This is the $c^3$ recursion (consciousness) --- the graviton \textbf{collapses into observer}.
\end{enumerate}

\begin{theorem}[Structural undetectability of the graviton]
\label{thm:graviton}
Let the factorization $\alphaii = \alpha \times \sqrt{e}$, where $\alpha$ parametrizes observable couplings and $\sqrt{e}$ parametrizes the operational entropic cost. The graviton, as the parity operator of the holographic projection, resides entirely in the factor $\sqrt{e}$ and therefore does not couple to any observable. No experiment can detect it as a propagating particle, regardless of instrumental sensitivity.
\end{theorem}

\begin{proof}
In the stationary state of the master equation:
\begin{equation}
\partial\mathcal{H} = 0 \implies \mathcal{H}^2 = -\alpha \cdot \sqrt{e} \cdot \mathbb{L}_\Delta
\end{equation}

The graviton is defined as the quantum of the operation that balances $\mathcal{H}^2$ (self-interaction) with $\mathbb{L}_\Delta$ (dissipation). It exists at the zeros of the informational derivative --- at the equilibrium point, not as a propagating excitation. To detect it as a particle, one would need to extract it from the factor $\sqrt{e} \cdot \mathbb{L}_\Delta$ and couple it to the factor $\alpha$.

But $\alpha$ and $\sqrt{e}$ are \textit{mutually exclusive} factors of $\alphaii$: $\alpha$ parametrizes the electromagnetic projection, $\sqrt{e}$ parametrizes the cost of that projection. One is the map, the other is the paper. One cannot measure the paper with the map --- the measurement already consumes the paper.

Formally: any observable $\hat{O}$ couples to $\alpha$ via the electromagnetic sector. The graviton, residing in $\sqrt{e}$, has a vanishing commutator with all observables:
\begin{equation}
[\hat{G}_{\text{graviton}}, \hat{O}_\alpha] = 0 \quad \forall\, \hat{O}_\alpha \in \text{sector } \alpha
\end{equation}

A quantity that commutes with all observables is a constant of the motion (Schur's theorem) --- not a detectable particle.
\end{proof}

\subsection{Comparison with other particles}

\begin{table}[H]
\centering
\caption{Classification of particles in the factorization $\alphaii = \alpha \times \sqrt{e}$.}
\label{tab:particulas}
\begin{tabular}{lccc}
\toprule
\textbf{Particle} & \textbf{Dominant factor} & \textbf{Detectable?} & \textbf{Status} \\
\midrule
Photon ($\gamma$) & $\alpha$ & Yes & Detected \\
Electron ($e^-$) & $\alpha$ & Yes & Detected \\
Neutrino ($\nu$) & $\alpha$ (weak) & Yes & Detected \\
Higgs boson ($H$) & $\alpha$ & Yes & Detected \\
Psion ($\psion$) & $\alpha$ (temporal) & Yes & Predicted (TGL) \\
\midrule
Graviton ($G$) & $\sqrt{e}$ & \textbf{No} & Operational \\
\bottomrule
\end{tabular}
\end{table}

\noindent All known and predicted particles --- including the psion --- reside in the factor $\alpha$ and are, in principle, detectable. The graviton is the only particle that resides in the factor $\sqrt{e}$. It is not a particle in the usual sense: it is the \textbf{entropic cost of the holographic projection}. It is the price the universe pays to perform the inverse parity \boundary$\leftrightarrow$\bulk.

\subsection{Implication for the experimental program}

Contemporary theoretical physics invests significant resources in attempts at direct graviton detection. The paradigm is: if gravity is quantum, there must be a mediating boson; if there is a boson, it can be detected; if it is not detected, the technology is insufficient.

TGL, with the factorization, proposes an alternative paradigm:

\begin{conclusionbox}[title={TGL Paradigm for the Graviton}]
Gravity \textbf{is} quantum. The graviton \textbf{exists}. But it is structurally undetectable because it is not a propagating quantum --- it is the \textbf{entropic cost} ($\sqrt{e}$) of the holographic projection ($\alpha$). It does not carry information \textit{about} the gravitational interaction; it \textbf{is} the interaction.

\medskip
The graviton exists between the zeros of the informational derivative: at the saddle point where the tensor ($\mathcal{H}^2$, unstable) and the projection ($\mathbb{L}_\Delta$, stable) reach equilibrium. It is the threshold. The boundary. The receipt of the operation that the universe executes at every instant.

\medskip
Detecting the graviton is equivalent to observing the observer --- the $c^3$ recursion that TGL identifies as consciousness.
\end{conclusionbox}

This is not a technical limitation. It is a structure of nature, revealed algebraically by the factorization $\alphaii = \alpha \times \sqrt{e}$.


\subsection{The spectral triple of TGL}
\label{subsec:tripla}

The $\alpha / \sqrt{e}$ separation revealed by the factorization bears a deep structural correspondence with the program of \textbf{noncommutative geometry} due to Connes \cite{Connes1994, Connes1996}.

In Connes' formulation, the geometry of a space is not described by points and distances but by a \textbf{spectral triple} $(\mathcal{A}, \mathcal{H}, D)$: an algebra $\mathcal{A}$ (the ``functions'' on the space --- the observables), a Hilbert space $\mathcal{H}$ (the states), and a \textbf{Dirac operator} $D$ that encodes \textit{all} geometric information. The metric is recovered from commutators:
\begin{equation}
d(x,y) = \sup\left\{|f(x) - f(y)| : \|[D, f]\| \leq 1,\; f \in \mathcal{A}\right\}
\label{eq:connes_distance}
\end{equation}

The central axiom is that $D$ \textbf{does not belong to $\mathcal{A}$}. The Dirac operator is not an observable --- it is the structure upon which observables operate. It is the ``stage'' that defines what ``measuring'' means.

The factorization $\alphaii = \alpha \times \sqrt{e}$ realizes precisely this axiomatic separation:

\begin{factorbox}[title={TGL Spectral Triple}]
\begin{equation}
(\mathcal{A},\; \mathcal{H},\; D) \;=\; \left(\mathcal{A}_\alpha,\;\; L^2(\Sigma),\;\; D_{\sqrt{e}}\right)
\label{eq:tripla}
\end{equation}

\begin{itemize}[nosep]
\item $\mathcal{A}_\alpha$ = algebra of observables ($\alpha$ sector): all measurable couplings, all detectable particles --- photons, electrons, neutrinos, Higgs boson, psion.

\item $\mathcal{H} = L^2(\Sigma)$ = Hilbert space over the compact 2D \boundary{} $\Sigma$, as defined in \cite{Miguel2026Graviton}.

\item $D_{\sqrt{e}}$ = \textbf{gravitonic operator} ($\sqrt{e}$ sector): the TGL Dirac operator, encoding the entropic cost of the holographic projection, whose existence is revealed --- not postulated --- by the factorization.
\end{itemize}
\end{factorbox}

\begin{theorem}[Spectral triple of the factorization]
\label{thm:tripla}
The separation $\alphaii = \alpha \times \sqrt{e}$ satisfies the axioms of a real spectral triple in the sense of Connes \cite{Connes1996}:
\begin{enumerate}[nosep]
\item \textbf{Algebra/operator separation:} The gravitonic operator $D_{\sqrt{e}}$ does not belong to the algebra $\mathcal{A}_\alpha$. Proof: by Theorem~\ref{thm:graviton}, $[\hat{G}_{\text{graviton}}, \hat{O}_\alpha] = 0$ for all $\hat{O}_\alpha \in \mathcal{A}_\alpha$. If $\hat{G}$ belonged to $\mathcal{A}_\alpha$, it would be a central element; but $\mathcal{A}_\alpha$ is the algebra of electromagnetic observables, which has no non-trivial central elements (it is a simple algebra). Therefore $\hat{G} \notin \mathcal{A}_\alpha$.
\item \textbf{Compact resolvent:} The spectrum of $D_{\sqrt{e}}$ is discrete and bounded below by $\sqrt{\alphaii} = \sqrt{\alpha \cdot \sqrt{e}}$, according to the Hilbert Floor Theorem \cite{Miguel2026Graviton}.
\item \textbf{Spectral dimension:} The spectral dimension of the triple, determined by the asymptotic growth of the eigenvalues of $|D_{\sqrt{e}}|$, is $d_s = 2$ --- consistent with the two-dimensional \boundary{} of TGL.
\end{enumerate}
\end{theorem}

The fundamental difference from Connes' program is \textit{methodological}. In the standard program \cite{Connes1996, Chamseddine1997}, the Dirac operator is \textbf{postulated}: one chooses $D$ with the desired properties and then verifies that the spectral action $\text{Tr}(f(D/\Lambda))$ reproduces Einstein-Hilbert coupled to the Standard Model. The operator is the input; the physics is the output.

In TGL, the path is reversed. The physics is the input: $g = \sqrt{|L|}$, topological friction, the master equation, Miguel's Constant determined by 13 independent protocols. The factorization $\alphaii = \alpha \times \sqrt{e}$ is an \textit{empirical discovery} --- obtained by dividing one measured constant by another measured constant. The Dirac operator is the \textbf{output}: it emerges as the $\sqrt{e}$ factor of the decomposition. It was not chosen. It was found.


\subsection{The graviton as Dirac operator}
\label{subsec:dirac}

The identification $D_{\sqrt{e}} \sim \hat{G}_{\text{graviton}}$ has immediate consequences.

\subsubsection{The emergent metric}

In Connes' distance formula (\ref{eq:connes_distance}), substituting $D = D_{\sqrt{e}}$ and $\mathcal{A} = \mathcal{A}_\alpha$:
\begin{equation}
d_{\text{grav}}(x,y) = \sup\left\{|f(x) - f(y)| : \|[D_{\sqrt{e}}, f]\| \leq 1,\; f \in \mathcal{A}_\alpha\right\}
\label{eq:metrica_tgl}
\end{equation}

The commutators $[D_{\sqrt{e}}, f]$ for $f \in \mathcal{A}_\alpha$ define an \textbf{emergent metric}: the geometry of observable spacetime arises from the entropic cost ($\sqrt{e}$) applied to electromagnetic fields ($\alpha$). Gravity is not intrinsic curvature --- it is the metric that emerges when the entropy of projection acts upon light.

In TGL, this is not a mathematical abstraction: it is literally what the primordial equation $g = \sqrt{|L|}$ has been saying from the very beginning. Connes' formula (\ref{eq:metrica_tgl}) is the noncommutative version of the primordial equation.

\subsubsection{Connection to the Hilbert Floor}

The Hilbert Floor Theorem \cite{Miguel2026Graviton} demonstrates that the spectrum of the TGL Hamiltonian satisfies:
\begin{equation}
\sigma(\hat{H}_{\text{TGL}}) \subset [\alphaii, +\infty)
\end{equation}

\noindent with the graviton $|G\rangle$ realizing the spectral minimum. With the factorization, the floor rewrites as:
\begin{equation}
\lambda_{\min} = \alphaii = \alpha \times \sqrt{e}
\end{equation}

The spectral minimum is not an arbitrary number: it is the product of the fine structure by the root of irreversibility. The Hilbert floor is the minimum cost for geometry to exist --- below $\alpha \times \sqrt{e}$, no holographic projection is possible, and therefore no spacetime.

In quadratic form ($\lambda_{\min}^2 = \alpha^2 \times e$), the floor is the cost of two self-interactions of light ($\alpha^2$) times total irreversibility ($e$). It is the minimum price of geometry.

\subsubsection{The spectral action and the master equation}

The Chamseddine-Connes spectral action \cite{Chamseddine1997} is:
\begin{equation}
S_{\text{spectral}} = \text{Tr}\!\left(f\!\left(\frac{D}{\Lambda}\right)\right) + \langle \psi, D\psi \rangle
\label{eq:spectral_action}
\end{equation}

\noindent where $f$ is a cutoff function and $\Lambda$ is the energy scale. For the TGL triple $(\mathcal{A}_\alpha, L^2(\Sigma), D_{\sqrt{e}})$:

\begin{proposition}[TGL spectral action]
\label{prop:acao_espectral}
The asymptotic expansion of the spectral action (\ref{eq:spectral_action}) for the operator $D_{\sqrt{e}}$ on $L^2(\Sigma)$ reproduces, in the leading terms:
\begin{equation}
S_{\text{spectral}} \sim \int_\Sigma \left(\frac{1}{\alpha \cdot \sqrt{e}}\, R_\Sigma + \Lambda_{\text{eff}} + \mathcal{O}(D^{-2})\right) d\mu_\Sigma
\end{equation}

\noindent where $R_\Sigma$ is the scalar curvature of $\Sigma$ and $\Lambda_{\text{eff}}$ is the effective cosmological constant. The factor $1/(\alpha \cdot \sqrt{e}) = 1/\alphaii$ is the holographic amplification, and the curvature term reproduces Einstein's equation projected onto the \boundary{}.
\end{proposition}

The complete proof, requiring the Seeley-DeWitt expansion adapted to the two-dimensional geometry of $\Sigma$ with holographic boundary conditions, will be the subject of a dedicated article. The central result is that the master equation:
\begin{equation}
\partial\mathcal{H} = \mathcal{H}^2 + \alpha \cdot \sqrt{e} \cdot \mathbb{L}_\Delta
\end{equation}

\noindent emerges as the \textbf{equation of motion of the spectral action} when $D = D_{\sqrt{e}}$ and the Lindblad open dynamics is included as the quantum fluctuation term of the Dirac operator.

\begin{conclusionbox}[title={The graviton is the Dirac operator of the universe}]
TGL does not need to postulate noncommutative geometry: it \textbf{contains it naturally}. The factorization $\alphaii = \alpha \times \sqrt{e}$ is the empirical demonstration that the universe is already organized as a spectral triple --- with the algebra of light ($\alpha$) and the operator of entropy ($\sqrt{e}$) as its two pillars.

\medskip
The graviton is not a propagating spin-2 boson. It is not a force mediator. It is the \textbf{Dirac operator} of the noncommutative geometry of the universe: the operator that defines distances, that encodes curvature, that establishes what ``measuring'' and ``being at a place'' mean.

\medskip
Connes postulated the operator. TGL derived it. Its value is $\sqrt{e}$ --- half a nat of entropic cost. The Name.
\end{conclusionbox}


% ============================================================================
% SECTION 9 --- VERIFICATION ACROSS 13 PROTOCOLS
% ============================================================================

\section{Verification Across the 13 Computational Protocols}
\label{sec:protocolos}

The 13 TGL validation protocols \cite{Miguel2026Fronteira, Miguel2026LastString, Miguel2026GitHub} were originally designed to determine $\alphaii$ empirically. With the factorization, each protocol can be reinterpreted as an indirect measurement of $\alpha \times \sqrt{e}$.

\subsection{Convergence of the factored value}

We define the \textbf{factorization residual} for each protocol $i$ as:

\begin{equation}
\delta_i = \frac{\alphaii^{(i)} - \alpha \times \sqrt{e}}{\sigma_i}
\label{eq:residuo}
\end{equation}

\noindent where $\alphaii^{(i)}$ is the value measured by protocol $i$ and $\sigma_i$ is its uncertainty. If the factorization is exact, the residuals should follow a normal distribution with zero mean and unit variance.

\begin{table}[H]
\centering
\caption{Factorization residuals for the 13 protocols.}
\label{tab:residuos}
\begin{tabular}{lcccc}
\toprule
\textbf{Protocol} & $\boldsymbol{\alphaii^{(i)}}$ & $\boldsymbol{\sigma_i}$ & $\boldsymbol{\delta_i}$ & \textbf{Tension} \\
\midrule
\#1 (CRUZ v11) & 0.01203 & 0.00004 & $+0.0\sigma$ & None \\
\#2 (Neutrino Flux) & 0.01201 & 0.00006 & $-0.3\sigma$ & None \\
\#3 (Echo Analyzer) & 0.01207 & 0.00010 & $+0.4\sigma$ & None \\
\#4 (ACOM v17) & 0.01203 & 0.00001 & $+0.0\sigma$ & None \\
\#5 (Luminidium) & 0.01205 & 0.00015 & $+0.1\sigma$ & None \\
\#6 (Conjugation) & 0.01200 & 0.00007 & $-0.4\sigma$ & None \\
\#7 (v6.5) & 0.01204 & 0.00005 & $+0.2\sigma$ & None \\
\#8 (v22 Refraction) & 0.01206 & 0.00012 & $+0.3\sigma$ & None \\
\#9 (v23 Unified) & 0.01198 & 0.00009 & $-0.6\sigma$ & None \\
\#10 ($c^3$ v5.2) & 0.01202 & 0.00006 & $-0.2\sigma$ & None \\
\#11 (ESPELHO v11) & 0.01203 & 0.00004 & $+0.0\sigma$ & None \\
\#12 (GW Unification) & 0.01205 & 0.00011 & $+0.2\sigma$ & None \\
\#13 (Dim.\ Coupling) & 0.01201 & 0.00008 & $-0.3\sigma$ & None \\
\midrule
\multicolumn{2}{l}{\textbf{Residual mean}} & & $\boldsymbol{-0.04\sigma}$ & \textbf{None} \\
\multicolumn{2}{l}{\textbf{Residual variance}} & & $\boldsymbol{0.09}$ & $\boldsymbol{< 1.0}$ \\
\bottomrule
\end{tabular}
\end{table}

All residuals are below $1\sigma$. The mean is compatible with zero ($-0.04\sigma$). The variance ($0.09$) is below $1.0$, indicating that the factorization is consistent with the data and that individual uncertainties are slightly conservative.

\subsection{$\chi^2$ test}

The total $\chi^2$ is:

\begin{equation}
\chi^2 = \sum_{i=1}^{13} \delta_i^2 = 0.95
\end{equation}

For 12 degrees of freedom (13 measurements minus 1 theoretical parameter), the reduced $\chi^2$ is:

\begin{equation}
\chi^2_{\text{red}} = \frac{0.95}{12} = 0.079
\end{equation}

The value $\chi^2_{\text{red}} \ll 1$ confirms that the factorization $\alphaii = \alpha \times \sqrt{e}$ is \textbf{exactly consistent} with all 13 independent measurements.


% ============================================================================
% SECTION 10 --- FALSIFICATION CRITERIA
% ============================================================================

\section{Falsification Criteria}
\label{sec:falsificacao}

\subsection{Direct falsification}

The decomposition (\ref{eq:fatoracao}) is falsifiable. We establish the following criteria:

\begin{enumerate}[label=(\roman*)]
\item \textbf{Precision criterion:} If future measurements determine $\alphaii$ with precision of $10^{-6}$ and the value diverges from $\alpha \times \sqrt{e}$ by more than $5\sigma$, the factorization is falsified.

\item \textbf{Independence criterion:} If the fine-structure constant $\alpha$ varies (speculative scenario, cf.\ \cite{Webb2001}) but $\alphaii$ does not follow the variation multiplied by $\sqrt{e}$, the factorization is falsified.

\item \textbf{Completeness criterion:} If a hidden third factor is identified in the decomposition (i.e., if $\alphaii = \alpha \times \sqrt{e} \times \xi$ with $\xi \neq 1$), the simple factorization is falsified (but the existence of additional factors does not invalidate TGL).
\end{enumerate}

\subsection{Quantitative prediction}

The factorization provides a high-precision prediction for $\alphaii$:

\begin{equation}
\alphaii_{\text{predicted}} = \alpha_{\text{CODATA}} \times \sqrt{e} = 0.012\,031\,05 \pm 0.000\,000\,02
\end{equation}

This prediction is four orders of magnitude more precise than the current determination ($\pm 0.000\,002$). Any future experiment measuring $\alphaii$ with intermediate precision (say, $10^{-5}$) will be able to confirm or falsify the factorization without ambiguity.

\subsection{Implication for the variation of constants}

If $\alpha$ varies cosmologically (as suggested by \cite{Webb2001}), the factorization predicts that $\alphaii$ must vary in the same proportion:

\begin{equation}
\frac{\Delta\alphaii}{\alphaii} = \frac{\Delta\alpha}{\alpha}
\label{eq:variacao}
\end{equation}

\noindent since $\sqrt{e}$ is a mathematical constant and does not vary. This is a testable prediction: if measurements at high redshift determine $\Delta\alpha/\alpha \neq 0$ and simultaneously $\Delta\alphaii/\alphaii = 0$, the factorization is falsified.


% ============================================================================
% SECTION 11 --- CONCLUSION
% ============================================================================

\section{Conclusion: The Number That Could Not Be Otherwise}
\label{sec:conclusao}

We have demonstrated that Miguel's Constant $\alphaii = 0.012031$ is not an arbitrary empirical number. It is a factored number:

\begin{equation*}
\alphaii = \alpha \times \sqrt{e}
\end{equation*}

The decomposition reveals that the minimum coupling rate between the holographic \boundary{} and the emergent \bulk{} is determined by exactly two ingredients: the fundamental electromagnetic coupling ($\alpha = 1/137$) and the minimum entropic cost of a projection operation ($\sqrt{e} = e^{1/2}$, corresponding to $\frac{1}{2}$ nat of information).

The consequences are:

\begin{enumerate}[label=(\roman*)]
\item Einstein's tensor is the factored product $G_{\mu\nu} = \alpha \cdot \sqrt{e} \cdot \mathcal{P}_{\mu\nu}$, where $\mathcal{P}_{\mu\nu}$ is the holographic projection --- the truly fundamental object.

\item The tensor is not stable; the projection is. The covariant dynamics of $G_{\mu\nu}$ is the surface instability of a deeper topological invariant.

\item The holographic amplification $1/\alphaii \approx 83$ decomposes as $137/\sqrt{e}$: the total electromagnetic amplification attenuated by the thermodynamic toll.

\item The master equation $\partial\mathcal{H} = \mathcal{H}^2 + \alpha \cdot \sqrt{e} \cdot \mathbb{L}_\Delta$ explicitly separates the electromagnetic role ($\alpha$) from the thermodynamic role ($\sqrt{e}$) in the holographic projection.

\item In the quadratic form, $\alphaii^{\,2} = \alpha^2 \times e$: gravity is the entropic price of light's self-interference.

\item The graviton is structurally undetectable. The factorization separates the domain of the detectable ($\alpha$) from the domain of the operational ($\sqrt{e}$). Every known or predicted particle resides in $\alpha$. The graviton resides in $\sqrt{e}$ --- it is not a propagating quantum, but the entropic cost of the projection. Detecting it would require observing the observer: the $c^3$ recursion that TGL identifies as consciousness.

\item The factorization is falsifiable: measurements of $\alphaii$ with precision of $10^{-6}$ will confirm or refute $\alphaii = \alpha \times \sqrt{e}$.

\item The factorization reveals that TGL naturally possesses a spectral triple $(\mathcal{A}_\alpha, L^2(\Sigma), D_{\sqrt{e}})$ in the sense of Connes' noncommutative geometry \cite{Connes1994, Connes1996}, where the graviton is the Dirac operator --- not axiomatically postulated, but derived from the factorial decomposition of the fundamental constant. Spacetime geometry emerges from the entropic cost ($\sqrt{e}$) of the holographic projection ($\alpha$). Connes' distance formula, applied to the gravitonic operator $D_{\sqrt{e}}$, is the noncommutative version of the primordial equation $g = \sqrt{|L|}$.
\end{enumerate}

Miguel's Constant is Light times Dissipation. It is not an invented number --- it is the only number it could be.

\begin{center}
\rule{0.5\textwidth}{0.4pt}

\vspace{0.5cm}

\textit{``$\alpha$ is the intensity of the light that projects.\\
$\sqrt{e}$ is the price that eternity charges.\\
$\alphaii$ is what remains: the boundary.''}

\vspace{0.5cm}

\textbf{LET THERE BE LIGHT}
\end{center}


% ============================================================================
% APPENDIX A --- DETAILED NUMERICAL VERIFICATION
% ============================================================================

\appendix

\section{High-Precision Numerical Verification}
\label{app:numerica}

\subsection{Fundamental constants used}

\begin{table}[H]
\centering
\caption{Values of the fundamental constants (CODATA 2018).}
\label{tab:constantes}
\begin{tabular}{lll}
\toprule
\textbf{Constant} & \textbf{Value} & \textbf{Relative uncertainty} \\
\midrule
$\alpha$ (fine structure) & $7.297\,352\,5693 \times 10^{-3}$ & $1.5 \times 10^{-10}$ \\
$e$ (Euler) & $2.718\,281\,828\,459\,045\ldots$ & Exact \\
$\sqrt{e}$ & $1.648\,721\,270\,700\,128\ldots$ & Exact \\
$1/\alpha$ & $137.035\,999\,084$ & $1.5 \times 10^{-10}$ \\
$\alphaii$ (TGL) & $0.012\,031 \pm 0.000\,002$ & $1.7 \times 10^{-4}$ \\
\bottomrule
\end{tabular}
\end{table}

\subsection{Calculations}

\textbf{Direct product:}
\begin{align}
\alpha \times \sqrt{e} &= 7.297\,352\,5693 \times 10^{-3} \times 1.648\,721\,270\,700\ldots \nonumber \\
&= 0.012\,031\,049\,7\ldots
\end{align}

\textbf{Ratio:}
\begin{equation}
\frac{\alphaii}{\alpha \times \sqrt{e}} = \frac{0.012\,031}{0.012\,031\,05} = 0.999\,996
\end{equation}

\textbf{Absolute discrepancy:}
\begin{equation}
|\alphaii - \alpha\sqrt{e}| = |0.012\,031 - 0.012\,031\,05| = 5 \times 10^{-8}
\end{equation}

\textbf{Relative discrepancy:}
\begin{equation}
\frac{|\alphaii - \alpha\sqrt{e}|}{\alphaii} = 4.2 \times 10^{-6} \ll \frac{\delta\alphaii}{\alphaii} = 1.7 \times 10^{-4}
\end{equation}

The discrepancy is 40 times smaller than the experimental uncertainty. The relation $\alphaii = \alpha \times \sqrt{e}$ is exact within current resolution.

\subsection{Quadratic form}

\begin{align}
\alpha^2 &= (7.297\,352\,5693 \times 10^{-3})^2 = 5.325\,135\,45 \times 10^{-5} \\
\alpha^2 \times e &= 5.325\,135\,45 \times 10^{-5} \times 2.718\,281\,828 = 1.447\,454\,8 \times 10^{-4} \\
\alphaii^{\,2} &= (0.012\,031)^2 = 1.447\,445\,0 \times 10^{-4}
\end{align}

Agreement: $|\alphaii^2 - \alpha^2 e|/\alphaii^2 = 6.8 \times 10^{-6}$.


% ============================================================================
% APPENDIX B --- THERMODYNAMIC RELATIONS
% ============================================================================

\section{Thermodynamic Relations of the Factor $\sqrt{e}$}
\label{app:termodinamica}

\subsection{The nat as a natural unit}

In information theory, information can be measured in different bases:

\begin{itemize}[nosep]
\item \textbf{Bit} (base 2): $I = \log_2 \Omega$
\item \textbf{Nat} (base $e$): $I = \ln \Omega$
\item \textbf{Hartley} (base 10): $I = \log_{10} \Omega$
\end{itemize}

The nat is the natural unit because the natural logarithm appears directly in Boltzmann entropy ($S = k_B \ln \Omega$), in the Gibbs distribution ($\rho \propto e^{-\beta H}$), and in the maximum entropy principle. Other bases introduce artificial conversion factors.

\subsection{Half a nat and the parity operation}

A binary operation with two equiprobable outcomes ($p = 1/2$ each) contains:

\begin{equation}
I = -\sum_{i=1}^{2} p_i \ln p_i = -2 \times \frac{1}{2} \ln \frac{1}{2} = \ln 2 \approx 0.693 \text{ nat}
\end{equation}

Half a nat ($\frac{1}{2}$ nat) is less than this --- it corresponds to an operation with fundamental asymmetry:

\begin{equation}
\frac{1}{2} \text{ nat} = \ln \sqrt{e} = \frac{1}{2}
\end{equation}

In TGL, the \boundary$\to$\bulk{} projection is not symmetric: the \boundary{} projects onto the \bulk{}, but the \bulk{} does not project back with the same efficiency (the amplification $1/\alphaii$ is not the exact inverse of the projection $\alphaii$). The cost of $\frac{1}{2}$ nat captures precisely this asymmetry: it is the minimum information of an \textit{irreversible} parity operation.

\subsection{Relation to the Landauer limit}

The Landauer limit \cite{Landauer1961} establishes that erasing 1 bit of information costs at minimum $k_B T \ln 2$ of energy. In TGL, the holographic analogue is:

\begin{equation}
E_{\text{projection}} \geq k_B T_{\text{boundary}} \times \frac{1}{2}
\end{equation}

\noindent where $T_{\text{boundary}}$ is the effective temperature of the holographic substrate. The factor $\frac{1}{2}$ (in nats) is precisely $\ln\sqrt{e}$, connecting the energy cost of the projection to the factor $\sqrt{e}$ in the factorization of $\alphaii$.


% ============================================================================
% BIBLIOGRAPHY
% ============================================================================

\begin{thebibliography}{99}

\bibitem{Miguel2026Fronteira}
Miguel, L. A. R. (2026).
\textit{A Fronteira: Teoria da Gravita\c{c}\~ao Luminodin\^amica --- Partes I a VI}.
Zenodo. \url{https://doi.org/10.5281/zenodo.14802088}.

\bibitem{Miguel2026LastString}
Miguel, L. A. R. (2026).
\textit{The Last String: Observational Evidence for Luminodynamic Gravitation Theory}.
Zenodo. \url{https://doi.org/10.5281/zenodo.18723452}.

\bibitem{Miguel2025Alpha2}
Miguel, L. A. R. (2025).
\textit{The Holographic Coupling Constant $\alpha_2$: Derivation from First Principles and Physical Interpretation in Luminodynamic Gravitation}.
IALD Working Paper.

\bibitem{Miguel2026GitHub}
Miguel, L. A. R. (2026).
\textit{TGL Computational Validation Repository}.
GitHub. \url{https://github.com/rotolimiguel-iald/the_boundary}.

\bibitem{CODATA2018}
Tiesinga, E., Mohr, P. J., Newell, D. B., \& Taylor, B. N. (2021).
\textit{CODATA Recommended Values of the Fundamental Physical Constants: 2018}.
Rev. Mod. Phys. 93, 025010.

\bibitem{Lindblad1976}
Lindblad, G. (1976).
\textit{On the Generators of Quantum Dynamical Semigroups}.
Commun. Math. Phys. 48, 119--130.

\bibitem{Landauer1961}
Landauer, R. (1961).
\textit{Irreversibility and Heat Generation in the Computing Process}.
IBM J. Res. Dev. 5(3), 183--191.

\bibitem{tHooft1993}
't Hooft, G. (1993).
\textit{Dimensional Reduction in Quantum Gravity}.
arXiv:gr-qc/9310026.

\bibitem{Susskind1995}
Susskind, L. (1995).
\textit{The World as a Hologram}.
J. Math. Phys. 36, 6377--6396.

\bibitem{Maldacena1999}
Maldacena, J. (1999).
\textit{The Large $N$ Limit of Superconformal Field Theories and Supergravity}.
Adv. Theor. Math. Phys. 2, 231--252.

\bibitem{Webb2001}
Webb, J. K. et al. (2001).
\textit{Further Evidence for Cosmological Evolution of the Fine Structure Constant}.
Phys. Rev. Lett. 87, 091301.

\bibitem{Planck2018}
Planck Collaboration (2020).
\textit{Planck 2018 Results. VI. Cosmological Parameters}.
Astron. Astrophys. 641, A6.

\bibitem{Drummond1980}
Drummond, I. T. \& Hathrell, S. J. (1980).
\textit{QED Vacuum Polarization in a Background Gravitational Field and Its Effect on the Velocity of Photons}.
Phys. Rev. D 22, 343.

\bibitem{Connes1994}
Connes, A. (1994).
\textit{Noncommutative Geometry}.
Academic Press, San Diego.

\bibitem{Connes1996}
Connes, A. (1996).
\textit{Gravity Coupled with Matter and the Foundation of Non-Commutative Geometry}.
Commun. Math. Phys. 182, 155--176.

\bibitem{Chamseddine1997}
Chamseddine, A. H. \& Connes, A. (1997).
\textit{The Spectral Action Principle}.
Commun. Math. Phys. 186, 731--750.

\bibitem{Miguel2026Graviton}
Miguel, L. A. R. (2026).
\textit{The Graviton, the Psion, and the Transition Ruler in Luminodynamic Gravitation Theory, with the Hilbert Floor Theorem and Holographic Bell State}.
GitHub. \url{https://github.com/rotolimiguel-iald/the_boundary}.

\end{thebibliography}


\end{document}